\chapter{QUANTUM MACHINE LEARNING (QML)}

This block focuses on the intersection between quantum computing and machine learning, where parametrized quantum circuits are used as trainable models.

\begin{center}
\textbf{Quantum Machine Learning combines quantum circuits with classical optimization to build models capable of learning from data}
\end{center}

\vspace{0.5cm}

\section*{1. FUNDAMENTAL IDEA}

Quantum Machine Learning relies on hybrid quantum-classical models:

\[
f(\theta, x) = \langle \psi(x, \theta) | M | \psi(x, \theta) \rangle
\]

where:
\begin{itemize}
\item $x$ is input data
\item $\theta$ are trainable parameters
\item $M$ is a measurement operator
\item $|\psi(x, \theta)\rangle$ is a parametrized quantum state
\end{itemize}

The learning process consists of minimizing a loss function:

\[
\min_\theta \mathcal{L}(f(\theta, x), y)
\]

\vspace{0.5cm}

\section*{2. QUANTUM NEURAL NETWORKS (QNNs)}

\subsection*{2.1 Definition}

Quantum Neural Networks are parametrized quantum circuits that act as function approximators.

\subsection*{2.2 Structure}

\begin{enumerate}
\item Encoding of classical data into quantum states
\item Application of parametrized gates
\item Measurement to extract outputs
\end{enumerate}

\subsection*{2.3 Mathematical Formulation}

\[
|\psi(x, \theta)\rangle = U(\theta) U(x) |0\rangle
\]

Output:

\[
y = \langle \psi(x, \theta) | M | \psi(x, \theta) \rangle
\]

\subsection*{2.4 Key Properties}

\begin{itemize}
\item Nonlinear function approximation via measurement
\item High-dimensional Hilbert space representation
\item Potential quantum advantage in specific tasks
\end{itemize}

\vspace{0.5cm}

\section*{3. VARIATIONAL QUANTUM EIGENSOLVER (VQE)}

\subsection*{3.1 Objective}

To approximate the ground state energy of a Hamiltonian:

\[
H |\psi\rangle = E |\psi\rangle
\]

\subsection*{3.2 Variational Principle}

\[
E(\theta) = \langle \psi(\theta) | H | \psi(\theta) \rangle
\]

\[
E(\theta) \geq E_0
\]

where $E_0$ is the true ground state energy.

\subsection*{3.3 Algorithm Structure}

\begin{enumerate}
\item Prepare parametrized state
\item Measure expectation value
\item Optimize parameters
\end{enumerate}

\subsection*{3.4 Applications}

\begin{itemize}
\item Quantum chemistry
\item Materials science
\item Ground state estimation
\end{itemize}

\vspace{0.5cm}

\section*{4. QUANTUM APPROXIMATE OPTIMIZATION ALGORITHM (QAOA)}

\subsection*{4.1 Objective}

To solve combinatorial optimization problems.

\subsection*{4.2 State Construction}

\[
|\gamma, \beta\rangle = e^{-i\beta H_M} e^{-i\gamma H_C} |+\rangle
\]

\subsection*{4.3 Optimization Problem}

\[
\max_{\gamma, \beta} \langle \gamma, \beta | H_C | \gamma, \beta \rangle
\]

\subsection*{4.4 Key Features}

\begin{itemize}
\item Alternating operator structure
\item Parameterized depth $p$
\item Suitable for NISQ devices
\end{itemize}

\vspace{0.5cm}

\section*{5. INTEGRATIONS}

\subsection*{5.1 Qiskit Machine Learning}

Provides tools for building hybrid models:

\begin{itemize}
\item Quantum kernels
\item Variational classifiers
\item Neural network interfaces
\end{itemize}

Example:

\begin{minted}{python}
from qiskit_machine_learning.neural_networks import EstimatorQNN
\end{minted}

\subsection*{5.2 PennyLane AI}

Enables seamless integration between quantum circuits and ML frameworks.

\begin{itemize}
\item Differentiable quantum programming
\item Integration with PyTorch, TensorFlow, JAX
\item Hybrid model training
\end{itemize}

Example:

\begin{minted}{python}

import pennylane as qml

dev = qml.device("default.qubit", wires=1)

@qml.qnode(dev)
def circuit(theta):
    qml.RY(theta, wires=0)
    return qml.expval(qml.PauliZ(0))
\end{minted}

\vspace{0.5cm}

\section*{6. TRAINING PROCESS}

The hybrid training loop:

\begin{enumerate}
\item Input data $x$
\item Prepare quantum state
\item Measure output
\item Compute loss
\item Update parameters using classical optimizer
\end{enumerate}

Mathematically:

\[
\theta_{k+1} = \theta_k - \eta \nabla_\theta \mathcal{L}
\]

\vspace{0.5cm}

\section*{FINAL CONNECTION}

Quantum Machine Learning unifies:

\begin{enumerate}
\item Quantum state representation
\item Parametrized circuit design
\item Classical optimization
\item Data-driven learning
\end{enumerate}

It extends classical machine learning into quantum Hilbert spaces.

\vspace{0.5cm}

\section*{STUDY ROADMAP}

\subsection*{Phase 1 --- Foundations}
\begin{itemize}
\item Parametrized quantum circuits
\item Measurement theory
\end{itemize}

\subsection*{Phase 2 --- Core Models}
\begin{itemize}
\item QNNs
\item VQE
\item QAOA
\end{itemize}

\subsection*{Phase 3 --- Integration}
\begin{itemize}
\item Qiskit Machine Learning
\item PennyLane AI
\end{itemize}

\subsection*{Phase 4 --- Training}
\begin{itemize}
\item Loss functions
\item Optimization methods
\end{itemize}

\subsection*{Phase 5 --- Applications}
\begin{itemize}
\item Classification
\item Optimization
\item Quantum simulation
\end{itemize}